\chapter*{Introduction generale}
\addcontentsline{toc}{chapter}{Introduction generale}
\markboth{Introduction generale}{Introduction generale}

Dans un systeme d'information moderne, le nombre de vulnerabilites remontees chaque
semaine par un scanner comme Qualys se compte en milliers. Chaque vulnerabilite est
associee a un score technique -- le QDS (Qualys Detection Score) -- qui mesure la
gravite intrinseque de la faille, independamment de l'actif sur lequel elle a ete
detectee. Or, un serveur de developpement isole du reseau et une base de donnees de
production accessible depuis Internet ne presentent pas le meme risque pour
l'entreprise, meme s'ils partagent exactement les memes vulnerabilites. Traiter
toutes les alertes avec la meme priorite revient soit a noyer les equipes de
remediation sous un volume ingerable, soit, a l'inverse, a laisser filer des risques
critiques faute de hierarchisation.

C'est ce probleme -- passer d'un inventaire brut de vulnerabilites a une
hierarchisation actionnable du risque -- que ce projet de fin d'annee cherche a
resoudre. Realise au sein de \organisme, il a donne naissance a une plateforme qui
croise deux sources de donnees historiquement traitees separement :
les exports de scan Qualys (la probabilite d'exploitation technique) et la CMDB (la
criticite metier et les mesures de protection deja en place sur chaque actif). De ce
croisement est derive un score unique et explicable par actif, l'\emph{Asset Risk
Score}, qui pilote ensuite des delais de remediation differencies (SLA) et permet de
mesurer objectivement la performance des equipes charge es de corriger les failles
(MTTR pondere).

Ce rapport documente l'ensemble de la demarche suivie. Le premier chapitre situe le
projet dans son contexte organisationnel et technique, formule la problematique et
delimite precisement son perimetre. Le deuxieme chapitre specifie les besoins
fonctionnels et non fonctionnels de la plateforme a travers une analyse UML des cas
d'utilisation. Le troisieme chapitre presente la conception logicielle retenue :
architecture globale, modele de donnees et diagrammes de sequence des principaux
flux. Le quatrieme chapitre, coeur scientifique du projet, formalise et demontre
mathematiquement le modele de risque -- ARS, classification SLA, MTTR pondere et
agregation en un Enterprise Risk Score. Le cinquieme chapitre decrit la realisation
technique de la plateforme, de l'ingestion des donnees a l'assistant conversationnel.
Le sixieme chapitre explique, graphique par graphique, comment lire le tableau de
bord produit par l'application. Le septieme chapitre presente la strategie de test et
valide numeriquement le modele sur un exemple complet. Le rapport se conclut par un
bilan critique du travail realise et par les perspectives d'evolution envisagees pour
la plateforme.

\clearpage
